\documentclass[11pt]{article}
\usepackage{amsmath,amssymb,amsthm}
\usepackage{booktabs}
\usepackage[margin=1in]{geometry}

\newtheorem{theorem}{Theorem}
\newtheorem{lemma}[theorem]{Lemma}
\newtheorem{corollary}[theorem]{Corollary}

\title{Problem 755: Not Zeckendorf}
\author{Project Euler Solutions}
\date{}

\begin{document}
\maketitle


\section{Problem Statement}
$f(n)$ = number of representations of $n$ as sum of distinct Fibonacci numbers (allowing consecutive). $S(n) = \sum_{{k=0}}^n f(k)$. Given $S(100)=415, S(10^4)=312807$. Find $S(10^{{13}})$.

\section{Mathematical Analysis}
Zeckendorf's theorem says every $n$ has a unique representation using non-consecutive Fibonacci numbers. Here we allow all subsets of Fibonacci numbers.

$f(n)$ counts the number of subsets $A$ of Fibonacci numbers with $\sum_{{a\in A}} a = n$. This is related to the binary representation in the Fibonacci numeral system, but with carries allowed.

The generating function is $\prod_{{k\ge 1}} (1 + x^{{F_k}})$. The sum $S(n)$ can be computed using a digit-DP on the Zeckendorf representation of $n$.

\section{Answer}

\[
\boxed{2877071595975576960}
\]

\end{document}
